\phantomsection\addcontentsline{toc}{section}{Biomedical Imaging}
\ECchapter{26}{Biomedical Imaging}{Turning physical signals into reliable clinical information}{ECRed}

\ECObjectives{
\begin{itemize}
\item Explain the physical principles and complete signal chain of major imaging techniques.
\item Analyse resolution, contrast, noise, artefacts and acquisition-time trade-offs.
\item Calculate echo depth and define validation and safety requirements for a clinical protocol.
\end{itemize}}

\begin{multicols}{2}
\begin{engineeringnews}
Biomedical imaging allows clinicians to inspect anatomy and function without surgery. Engineers combine wave
physics, detectors, analogue electronics, digital signal processing and inverse algorithms to obtain clinically
useful information while controlling acquisition time, cost, uncertainty and patient risk.
\end{engineeringnews}

\begin{ecarticle}
Biomedical imaging converts interactions inside the body into spatial maps. X-ray radiography and computed
tomography measure the attenuation of ionising radiation. Ultrasound transmits mechanical waves and analyses
their echoes. Magnetic resonance imaging, or MRI, detects radio-frequency signals produced by hydrogen nuclei
in a strong magnetic field. Nuclear medicine detects radiation emitted by a tracer introduced into the body.

Every imaging system contains four linked functions: excitation, interaction, detection and reconstruction.
The detector does not usually deliver a ready-to-read image. It records many incomplete and noisy measurements
from different positions or times. A reconstruction algorithm estimates the spatial distribution that is most
consistent with these measurements and with the physical model of the system.

Image quality cannot be reduced to one number. Spatial resolution determines the smallest visible detail,
contrast separates tissues with similar properties, and signal-to-noise ratio indicates whether a weak feature
can be distinguished from random fluctuations. Artefacts are structured errors caused by motion, metal,
calibration defects, limited viewing angles or an incorrect model. Improving one criterion often degrades
another: smaller voxels may increase acquisition time or radiation dose, while faster scans may reduce signal.

In pulse-echo ultrasound, reflector depth is estimated from
\[
d=\frac{ct}{2},
\]
where $c$ is the speed of sound in tissue and $t$ is the round-trip delay. The factor two accounts for the outward
and return paths. Axial resolution depends on pulse duration, whereas lateral resolution depends on beam width
and focusing. A protocol must therefore specify both geometry and signal-processing parameters.

Safety constraints depend on the technology. X-ray dose must be kept as low as reasonably achievable. MRI
avoids ionising radiation but requires strict control of ferromagnetic objects, radio-frequency heating and
acoustic noise. Ultrasound is portable and comparatively inexpensive, yet image quality depends strongly on
operator technique and acoustic access.

Artificial intelligence can assist segmentation, reconstruction and anomaly detection, but it must be validated
on diverse patients, scanners and operating conditions. A plausible-looking image is not automatically a faithful
representation of the measurements. Engineers must quantify uncertainty, failure modes and traceability before
clinical deployment.
\end{ecarticle}

\begin{tcolorbox}[title=\sffamily\bfseries Engineering insight,colback=yellow!10,colframe=ECOrange]
A reconstructed image is an estimate, not a direct photograph. Its reliability depends on the measurement model,
calibration, sampling geometry and regularisation assumptions.
\end{tcolorbox}

\begin{engineeringfocus}
\textbf{From physical interaction to clinical decision}
\begin{center}
\resizebox{0.96\linewidth}{!}{%
\begin{tikzpicture}[font=\scriptsize,>=Latex,node distance=3.5mm]
\node[draw=ECRed,fill=ECRed!10,rounded corners] (source) {excitation};
\node[draw=ECOrange,fill=ECOrange!10,rounded corners,right=of source] (body) {body interaction};
\node[draw=ECBlue,fill=ECLightBlue,rounded corners,right=of body] (det) {detector};
\node[draw=ECPurple,fill=ECPurple!10,rounded corners,below=of body] (alg) {reconstruction};
\node[draw=ECGreen,fill=ECGreen!10,rounded corners,below=of det] (image) {validated image};
\draw[->,thick] (source)--(body); \draw[->,thick] (body)--(det);
\draw[->,thick] (det)--(alg); \draw[->,thick] (alg)--(image);
\draw[->,thick,dashed] (image) to[bend left=25] node[below]{quality feedback} (det);
\end{tikzpicture}}
\end{center}
\end{engineeringfocus}

\begin{keyfigures}
\begin{tabularx}{\linewidth}{@{}Xr@{}}
X-ray / CT & attenuation\\
Ultrasound & echo delay\\
MRI & magnetic resonance\\
Main trade-off & quality--time--risk\\
Critical requirement & clinical validation
\end{tabularx}
\end{keyfigures}

\begin{tcolorbox}[title=\sffamily\bfseries Resolution and acquisition time,colback=white,colframe=ECRed]
\begin{center}
\begin{tikzpicture}
\begin{axis}[width=0.94\linewidth,height=45mm,xlabel={Voxel size (mm)},ylabel={Acquisition time (s)},
grid=major,ticklabel style={font=\scriptsize},label style={font=\scriptsize},x dir=reverse]
\addplot[thick,mark=*] table[x=resolution,y=time,col sep=comma]{data/EC26/imaging.csv};
\end{axis}
\end{tikzpicture}
\end{center}
\footnotesize Smaller voxels may improve geometric detail, but only if the signal-to-noise ratio remains sufficient.
\end{tcolorbox}

\begin{applicationexercise}{Ultrasound depth and uncertainty}
An echo returns from a tissue boundary after $130\,\mu\mathrm{s}$. Take $c=1540\,\mathrm{m\,s^{-1}}$.
\begin{enumerate}
\item Calculate the depth of the boundary.
\item A second echo arrives $26\,\mu\mathrm{s}$ later. Estimate the distance between the two boundaries.
\item If the assumed sound speed is uncertain by $\pm 20\,\mathrm{m\,s^{-1}}$, estimate the resulting uncertainty on the first depth.
\item Explain why the full travel distance would give a physically incorrect result.
\end{enumerate}
\end{applicationexercise}

\begin{vocabulary}
\begin{tabularx}{\linewidth}{@{}lX@{}}
attenuation & atténuation\\
echo & écho\\
voxel & voxel\\
spatial resolution & résolution spatiale\\
contrast & contraste\\
noise & bruit\\
artefact & artefact\\
reconstruction & reconstruction\\
radiation dose & dose de rayonnement\\
traceability & traçabilité
\end{tabularx}
\end{vocabulary}

\begin{questions}
\item Why do different imaging technologies reveal different tissue properties?
\item What assumptions are introduced by reconstruction algorithms?
\item Distinguish random noise from a structured artefact.
\item Why is the highest possible spatial resolution not always the best protocol?
\item Give one safety constraint for CT and one for MRI.
\end{questions}

\begin{engineeringchallenge}
Design an imaging protocol for a moving organ. Define the clinical requirement, choose a technology, specify
spatial and temporal resolution, identify the dominant artefact, propose one correction method and build a
validation matrix containing at least three tests, one acceptance threshold and one degraded operating case.
\end{engineeringchallenge}

\begin{speaking}
Should AI-enhanced medical images be accepted when visibility improves but the original measurements are
modified? Discuss accuracy, traceability, responsibility and patient safety.
\end{speaking}
\end{multicols}

\subsection*{Sources}
\ECsource{World Health Organization -- Medical imaging}{https://www.who.int/health-topics/medical-devices}
\ECsource{International Atomic Energy Agency -- Radiation protection}{https://www.iaea.org/resources/rpop}